Los resultados obtenidos muestran que la corriente de ánodo presenta un incremento
monótono con el voltaje aplicado, comportamiento esperado para un diodo de vacío
operando bajo emisión termoiónica. Al realizar el ajuste de la corriente en función
de $V_a^{3/2}$ se obtuvo un coeficiente de determinación $R^2=0.995$, indicando que
la ley de Child--Langmuir describe adecuadamente los datos experimentales dentro del
intervalo analizado.\\

A partir de la pendiente del ajuste se calculó una relación carga--masa del electrón

\[
\frac{e}{m_e}=(1.719\pm0.032)\times10^{11}\ \mathrm{C/kg},
\]

valor que difiere únicamente un $2.23\%$ del valor aceptado
$\frac{e}{m_e}=1.7588\times10^{11}\ \mathrm{C/kg}$.
Esta diferencia se encuentra dentro de lo esperado para un montaje experimental de
laboratorio y sugiere que el procedimiento empleado permite estimar correctamente
esta constante fundamental.\\

El ajuste mediante una ley de potencia produjo un exponente
$p=1.243\pm0.013$, inferior al valor teórico $p=3/2$. Esta discrepancia indica que,
aunque la dependencia global entre corriente y voltaje sigue la tendencia prevista,
los datos experimentales no corresponden completamente al régimen ideal limitado por
carga espacial. Entre las posibles causas se encuentran fluctuaciones en la
temperatura del filamento, incertidumbre en la lectura de la corriente para voltajes
pequeños, variaciones en las dimensiones reales del tubo respecto a los valores
nominales y efectos asociados a la corriente de saturación.\\

En conjunto, los resultados muestran que la ley de Child--Langmuir constituye una
buena aproximación para describir el comportamiento del diodo termoiónico, aunque
las limitaciones experimentales afectan la determinación del exponente de potencia.\\

En conclusión, la práctica permitió comprobar experimentalmente el comportamiento
de la emisión termoiónica y estimar satisfactoriamente la relación carga--masa del
electrón utilizando un método alternativo al de los campos eléctricos y magnéticos.